\input{../tex-inputs/latex-leseansicht-vorspann}

\section[Gustav Schwarzkopf an Arthur Schnitzler, 24. 7. 1895]{L04282 Gustav Schwarzkopf an Arthur Schnitzler, 24. 7. 1895}
\nopagebreak\mylabel{L04282v}
\rehead{ }\normalsize\beginnumbering\briefempfaengerindex{Schnitzler, Arthur@\textsc{Schnitzler, Arthur}!zzzSchwarzkopf, Gustav@\emph{von Gustav Schwarzkopf}!1895-07-241@{24. 7. 1895}|(be}
\toendnotes[C]{\smallbreak\pagebreak[2]}
\correspDesc{Versand  durch Gustav Schwarzkopf am 24. 7. 1895 in Wien
\newline{}Erhalt  durch Arthur Schnitzler im Zeitraum [25. 7. 1895 – 29. 7. 1895?] in Wien}\toendnotes[C]{\smallbreak}
\Standort{CUL, Schnitzler, B 96a.}
\physDesc{Brief, 1 Blatt, 4 Seiten, 1716 Zeichen
\newline{}Handschrift: schwarze Tinte, deutsche Kurrent}
\pstart
           \raggedleft{}{\pb}\textsc{Wien\oindex{Wien@\textbf{Wien}, \emph{Verwaltungsgebiet}|pw}}{ }24. VII 95\pend
           
\pstart{}Lieber Freund!\pend\vspace{0.5em}
\pstart
           Sie{ }ſchildern \textsc{Ischl\oindex{Bad Ischl@\textbf{Bad Ischl}|pw}} mit allen{ }ſeinen Reizen ja verlockend genug, aber es iſt doch nicht
               wahrſcheinlich, daß ich die{ }ſtattliche Zahl{ }ſeiner Gäſte vermehren werde. Ich bin
               eben{ }ſo gar nicht in der Sti{\geminationm}ung, die für \textsc{Ischl\oindex{Bad Ischl@\textbf{Bad Ischl}|pw}} erforderlich iſt und ich würde nur Anderen die Laune verderben. In Wien\oindex{Wien@\textbf{Wien}, \emph{Verwaltungsgebiet}|pw} iſt es übrigens bis jetzt ganz erträglich und
                  we{\geminationn} ich ein Poet oder {\pb}eine Künſtlernatur wäre,{ }ſo kö{\geminationn}te ich auf meinen Spaziergängen überRing-\oindex{Wien@\textbf{Wien}!I., Innere Stadt@\textbf{I., Innere Stadt}!Ringstraße@\textbf{Ringstraße}, \emph{Straße}|pw} und Gürtelstraße\oindex{Gürtel@\textbf{Gürtel}, \emph{Straße}|pw} in
               dem{ }ſo{\geminationm}erlichen Wien\oindex{Wien@\textbf{Wien}, \emph{Verwaltungsgebiet}|pw}
               die{ }ſchönſten Entdeckungen machen; ich kö{\geminationn}te die
               träumende Großſtadtſeele belauſchen, die{ }ſich über sich{ }ſelbſt beugt, um das Leben zu
               deuten und dem ironiſch lächelnden Schickſal{ }ſeine letzten perverſen Heimlichkeiten
               abzuliſten – nein, es geht wirklich nicht, ich ka{\geminationn} nun
               einmal in dem Genre nichts leiſten. – Es iſt auch noch gar nicht{ }ſo leer hier, wie
               die abweſenden {\pb}oberen Zehntauſend
               protzig anzunehmen belieben, man trifft sogar noch Dichter, we{\geminationn} man Glück hat; vor wenigen, Tagen habe ich im Verlauf
               einer Viertelſtunde zwei geſehen: die Herren \textsc{Eber-\pwindex{Ebermann, Leo 16.\,7.\,1863 Draganovka – 9.\,10.\,1914 Wien@\textsc{Ebermann, Leo} (16.\,7.\,1863 Draganovka – 9.\,10.\,1914 Wien), \emph{Schriftsteller, Journalist, Rechtswissenschaftler}|pw}} und \textsc{Dörma{\geminationn}\pwindex{Dörmann, Felix 29.\,5.\,1870 Wien – 26.\,10.\,1928 ebd.@\textsc{Dörmann, Felix} (29.\,5.\,1870 Wien – 26.\,10.\,1928 ebd.), \emph{Schriftsteller}|pw}}.\pend
           
\pstart
           – In den erſten Tagen des Auguſt werde ich aber wahrſcheinlich doch fortgehen,
               nicht aus eigenem Antrieb, nicht um mich zu erholen, nur um der Schande zu entri{\geminationn}en, den ganzen Sommer hier geweſen zu{ }ſein, um ein \textsc{Alibi} nachweiſen zu kö{\geminationn}en, um
               im {\pb}Herbst auf die zahlreichen Fragen
               »Wo waren Sie?« doch etwas antworten zu kö{\geminationn}en. Wohin,
               weiß ich noch nicht und{ }ſicher iſt’s auch noch nicht. – Max\pwindex{Schwarzkopf, Max 12.\,6.\,1857 Wien – 14.\,4.\,1928 ebd.@\textsc{Schwarzkopf, Max} (12.\,6.\,1857 Wien – 14.\,4.\,1928 ebd.), \emph{Rechtsanwalt}|pw} iſt{ }ſchon seit 14 Tagen in \textsc{Sulden\oindex{Solda@\textbf{Solda}|pw}}. Haben Sie{ }ſchon von einem Autor Hans von \textsc{Kahlenberg}\pwindex{Keßler, Helene 23.\,2.\,1870 Heilbad Heiligenstadt – 8.\,8.\,1957 Baden-Baden@\textsc{Keßler, Helene} (23.\,2.\,1870 Heilbad Heiligenstadt – 8.\,8.\,1957 Baden-Baden), \emph{Schriftstellerin}|pw} etwas gehört? Man hat mir vor einigen Tagen ein Buch »Ein Narr\pwindex{Keßler, Helene 23.\,2.\,1870 Heilbad Heiligenstadt – 8.\,8.\,1957 Baden-Baden@\textsc{Keßler, Helene} (23.\,2.\,1870 Heilbad Heiligenstadt – 8.\,8.\,1957 Baden-Baden), \emph{Schriftstellerin}!Narr. Roman@\strich\emph{Ein Narr. Roman}|pw}« zugeſchickt. Wieder Einer, der Talent zu haben{ }ſcheint. Es iſt entſetzlich!\pend
           
\pstart
           Herzlichſt{\\[\baselineskip]}Ihr{\\[\baselineskip]}\spacefill\mbox{GustavSchwkpf}\pend
           \leftskip=0em{}
\pstart
           \noindent{}Bitte, grüßen Sie \textsc{Beer-Hofma{\geminationn}\pwindex{Beer-Hofmann, Richard 11.\,7.\,1866 Wien – 26.\,9.\,1945 New York City@\textsc{Beer-Hofmann, Richard} (11.\,7.\,1866 Wien – 26.\,9.\,1945 New York City), \emph{Schriftsteller}|pw}}\pend
           \selectlanguage{ngerman}\endnumbering\briefempfaengerindex{Schnitzler, Arthur@\textsc{Schnitzler, Arthur}!zzzSchwarzkopf, Gustav@\emph{von Gustav Schwarzkopf}!1895-07-241@{24. 7. 1895}|)be}\mylabel{L04282h}
\begin{anhang}
\end{anhang}\newcommand{\dateiname}{L04282}\newcommand{\titel}{Gustav Schwarzkopf an Arthur Schnitzler, 24. 7. 1895}\newcommand{\editorInnen}{Herausgegeben von Jahnke, SelmaMüller, Martin Anton}\input{../tex-inputs/latex-leseansicht-abspann}
